% Part: computability
% Chapter: recursive-functions
% Section: bounded-minimization

\documentclass[../../../include/open-logic-section]{subfiles}

\begin{document}

\olfileid{cmp}{rec}{bmi}
\olsection{کمینه‌سازی کراندار}

\begin{explain}
اغلب سودمند است تابعی را به‌عنوان کوچک‌ترین عددی تعریف کنیم که ویژگی
یا رابطه‌ای چون~$P$ را برآورده می‌کند. اگر $P$ تصمیم‌پذیر باشد،
می‌توانیم این تابع را به‌سادگی با آزمودن همهٔ اعداد ممکن،
$0$، $1$، $2$، \dots، محاسبه کنیم تا به کوچک‌ترین عددی برسیم که~$P$
را برآورده می‌کند. این گونه جست‌وجوی نامحدود ما را از قلمرو توابع
بازگشتی اولیه بیرون می‌برد. اما اگر تنها به کوچک‌ترین عددی علاقه‌مند
باشیم که 
\emph{از کرانی مستقل و ازپیش‌داده‌شده کوچک‌تر است}، همچنان در قلمرو
بازگشتی اولیه می‌مانیم. به بیانی اندکی کلی‌تر، فرض کنید رابطهٔ
بازگشتی اولیهٔ~$R(x,z)$ را داریم. تابعی را در نظر بگیرید که $x$ و~$y$
را به کوچک‌ترین $z<y$ می‌نگارد که $R(x,z)$. این تابع نیز با آزمودن
$R(x,0)$، $R(x,1)$، \dots، $R(x,y-1)$ محاسبه‌پذیر است. اما چرا
بازگشتی اولیه است؟
\end{explain}

\begin{prop}
اگر $R(\vec x,z)$ بازگشتی اولیه باشد، تابع~$m_R(\vec x,y)$ نیز
بازگشتی اولیه است؛ این تابع، اگر $z<y$ای وجود داشته باشد که
$R(\vec x,z)$، کوچک‌ترین چنین~$z$ای را بازمی‌گرداند و در غیر این صورت
$y$ را بازمی‌گرداند. تابع~$m_R$ را چنین می‌نویسیم:
\[
\bmin{z < y}{R(\vec{x}, z)}.
\]
\end{prop}

\begin{proof}
هیچ $z<0$ای نمی‌تواند وجود داشته باشد که $R(\vec x,z)$، زیرا اصلاً
هیچ $z<0$ای وجود ندارد. پس $m_R(\vec x,0)=0$.

اگر کران به صورت $y+1$ باشد، سه حالت داریم:
\begin{enumerate}
\item $z<y$ای وجود دارد چنان‌که $R(\vec{x},z)$؛ در این صورت
$m_R(\vec{x},y+1)=m_R(\vec{x},y)$.
\item چنین $z<y$ای وجود ندارد اما $R(\vec{x},y)$؛ در این صورت
$m_R(\vec{x},y+1)=y$.
\item هیچ $z<y+1$ای وجود ندارد چنان‌که $R(\vec{x},z)$؛ در این صورت
$m_R(\vec{x},y+1)=y+1$.
\end{enumerate}
پس می‌توانیم $m_R$ را با بازگشت اولیه چنین تعریف کنیم:
\begin{align*}
m_R(\vec{x}, 0) & = 0\\
m_R(\vec{x}, y+1) & =
\begin{cases}
m_R(\vec{x}, y) & \text{اگر $m_R(\vec x, y) \neq y$}\\
y & \text{اگر $m_R(\vec x, y) = y$ و $R(\vec{x}, y)$}\\
y+1 & \text{در غیر این صورت}.
\end{cases}
\end{align*}
توجه کنید که $z<y$ای با ویژگی~$R(\vec x,z)$ وجود دارد هرگاه و تنها
هرگاه $m_R(\vec x,y) \neq y$.
\end{proof}

\begin{prob}
فرض کنید $R(\vec x,z)$ بازگشتی اولیه باشد. تابع~$m'_R(\vec x,y)$ را
که، اگر $z<y$ای با ویژگی~$R(\vec x,z)$ وجود داشته باشد، کوچک‌ترین
چنین~$z$ای و در غیر این صورت~$0$ را بازمی‌گرداند، با بازگشت اولیه
از~$\Char{R}$ تعریف کنید.
\end{prob}

\end{document}
